\par\medskip
\noindent
\underline{5. El teorema de ortogonalidad en el caso $\ell=p$} (\underline{caso
semiestable}). \underline{Dualidad de los esquemas abelianos $B_o$ y $B'_o$.
Caracterización de la parte fija. Criterios de buena reducción,
reducción esencialmente buena y reducción semiestable.}

\medskip
\noindent
5.1. Situémonos en el contexto de 2.2.3, donde $S$ es henseliano y
$A^o$ es $\ell$-divisible, de modo que la \underline{parte fija}
\[
 T_\ell(A^o)^f\lhook\joinrel\longrightarrow T_\ell(A)
\]
está definida como el pro-$\ell$-grupo de Barsotti--Tate formado por
los $({}_{\ell^\nu}A^o)^f$. Daremos ahora una definición de la parte
tórica (2.3.2) de $T_\ell(A^o)$ independiente de la hipótesis
$\ell\neq p$ formulada en 2.3, pero que supone en su lugar que $S$ es
\underline{completo}.\footnote{Esta hipótesis no es en realidad
esencial, como se muestra en 6.1 más adelante.} Para ello, para todo
entero $n\geq0$, introduzcamos el esquema artiniano local
\begin{equation}\tag{5.1.1}
 S_n=\operatorname{Spec}(V/\mathfrak m^{n+1}),
\end{equation}
el $n$-ésimo entorno infinitesimal de $s$ en $S$, y, para todo
$S$-esquema $X$, pongamos $X_n=X\times_S S_n$. Puesto que
$(A^o)_n=(A_n)^o$ es liso sobre $S_n$, se sabe (SGA 3 IX 3.6 bis) que
$T_o$ es la fibra especial de un único subtoro
\begin{equation}\tag{5.1.2}
 T_n\subset A_n^o.
\end{equation}
Por la unicidad de los $T_n$, estos se obtienen unos de otros mediante
los cambios de base $S_{n'}\longrightarrow S_n$. En otras palabras,
definen un \underline{subtoro formal} $\widehat T$ del grupo
$\widehat A^o$ (en la categoría de esquemas formales sobre $S$), el
completado formal de $A^o$ a lo largo de la fibra especial:
\begin{equation}\tag{5.1.3}
 \widehat T\subset\widehat A^o.
\end{equation}

Puesto que el funtor $X\longmapsto\widehat X$ (completado formal a lo
largo de la fibra especial) induce una equivalencia entre la categoría
de los esquemas finitos sobre $S$ y la de los esquemas formales finitos
sobre $S$ (EGA III 4.8), identificaremos simplemente un esquema finito
sobre $S$ con el esquema formal correspondiente. Con esta
identificación, la definición de las \emph{partes fijas} da de inmediato,
para todo $\nu\geq0$,
\begin{equation}\tag{5.1.4}
 ({}_{\ell^\nu}A^o)^f={}_{\ell^\nu}(\widehat A^o),
\end{equation}
lo que puede escribirse también como una identidad de sistemas
proyectivos:
\begin{equation}\tag{5.1.5}
 T_\ell(A^o)^f=T_\ell(\widehat A^o).
\end{equation}
Obtenemos, por tanto, un prosubgrupo de Barsotti--Tate de la parte fija
de $T_\ell(A^o)$ tomando
$T_\ell(\widehat T)\subset T_\ell(\widehat A^o)$. Definimos en
consecuencia la \underline{parte tórica} de $T_\ell(A^o)$ por
\begin{equation}\tag{5.1.6}
 T_\ell(A^o)^t
 \stackrel{\mathrm{def}}{=}
 T_\ell(\widehat T)
 \subset T_\ell(\widehat A^o)
 =T_\ell(A^o)^f.
\end{equation}
Así, por definición,
\begin{equation}\tag{5.1.7}
 T_\ell(A^o)^t=({}_{\ell^\nu}\widehat T)_{\nu\geq0},
\end{equation}
describiéndose el propio $S$-grupo finito plano
${}_{\ell^\nu}\widehat T$ como el límite inductivo de los $S_n$-grupos
finitos planos
\begin{equation}\tag{5.1.8}
 {}_{\ell^\nu}(T_n),\qquad n\ \text{variable};
\end{equation}
es decir, su álgebra afín se describe como el límite proyectivo de las
álgebras afines de los grupos (5.1.8). Naturalmente,
\begin{equation}\tag{5.1.9}
 \bigl(T_\ell(A^o)^t\bigr)_n=T_\ell(T_n).
\end{equation}
Para $n=0$, esta es precisamente la fórmula (2.3.3), mediante la cual
caracterizamos la parte tórica (como grupo profinito étale) en el caso
$\ell\neq p$. Esto muestra que la definición (5.1.6) es compatible con
la dada en 2.3 cuando $\ell\neq p$.

Pasando a las fibras genéricas, las partes fija y tórica de
$T_\ell(A^o)$ definen partes fija y tórica en $T_\ell(A_K)$:
\begin{equation}\tag{5.1.10}
 T_\ell(A_K)^t\subset T_\ell(A_K)^f\subset T_\ell(A_K).
\end{equation}
Estas son inclusiones de pro-$\ell$-grupos de Barsotti--Tate.
Considérese la filtración análoga correspondiente al esquema abeliano
dual $A'_K$:
\begin{equation}\tag{5.1.11}
 T_\ell(A'_K)^t\subset T_\ell(A'_K)^f\subset T_\ell(A'_K).
\end{equation}
Tenemos entonces el resultado siguiente, que generaliza 2.4 (al menos
en el caso en que $S$ es completo, al cual podríamos habernos reducido
fácilmente en 2.4).
